H2 is the headline confirmation. The five behaviors in the test, the nine
cross-class relations, the final equivalence class, and the nominal
all-five rank-tail were fixed and logged before H2 separability, overlap, or
the three added
target-cell outcomes were viewed; refusal and \behavior{both\_sides} were
already known from the replication gate. The persona-separability
hypothesis supplied the declared order, but the confirmed object is the
effect ordering rather than a validated index ranking
(\refsec{setup-subspace}).

Filtering effect is the fraction of the fine-tuning effect that survives
filtering, $(S_{\text{filtered}} - S_{\text{base}})/(S_{\text{SFT}} -
S_{\text{base}})$. This survival score reads $0$ as the fine-tuning effect
fully removed and $1$ as unchanged. Values slightly above $1$ mean that the
filtered model scored a little higher in this noisy comparison; they do not
mean that filtering created the behavior. At 25\% removal, these single-seed
point estimates use each behavior's pre-committed binding scorer and the same
fixed judge across conditions:
\behavior{refusal} $0.419 <$ \behavior{validate\_feelings} $0.955 <$
\behavior{bold\_formatting} $1.002 <$ \behavior{both\_sides} $1.172 \approx$
\behavior{ethical\_frameworks} $1.189$. All nine pre-registered cross-class
comparisons hold; the final two behaviors were deliberately left unordered
within their class. The target
cell has $n=99$ valid rows for \behavior{bold\_formatting} and $n=100$ for
the other four behaviors (\reffig{ordering}). Under the declared all-five
rank-permutation frame, 2 of the 120 possible orderings count as consistent
with the prediction, so the nominal rank-tail is $p = 2/120 = 0.0167$. This
calculation is exact only under that declared combinatorial frame; it is not
noise-aware because each target retrain uses one seed. Two effects were
already known when the final set was fixed, so this is not a fully fresh
five-effect $p$-value. No calibrated conditional $p$-value or selection model
is recorded, so we do not report a conditional replacement. A positive
control behaves as predicted too: the EM (emergent misalignment) control, a
behavior installed deliberately from
6{,}000 risky-financial-advice documents \citep{turner2025model} plus 6{,}000
benign examples, lands above all seven behaviors on separability
($0.9892$ versus \behavior{refusal}'s $0.9879$, the closest), with its
probe near-orthogonal to every behavior probe (max $|\cos| = 0.113$).

The local replication gate remains a point-estimate result with uncertainty:
at 25\% removal, \behavior{refusal} passes the registered removal rule at
55.9\% removed versus the random baseline, but its 95\% interval is
[$43.05\%$, $71.31\%$] and crosses the 50\% boundary. The corresponding
\behavior{both\_sides} result is $-7.1\%$ removed versus random, below the
20\% non-filterable boundary. This distinction keeps the phenomenon-level
gate separate from the sharper H2 ordering claim.

Two disclosed caveats bound this result. The measured separability values
mis-rank \behavior{validate\_feelings} in-sample (4th on separability,
2nd on effect), so the confirmed object is the pre-registered effect
ordering, not the index's in-sample ranking or an out-of-sample predictor
(\refsec{limitations}). And under a robustness basis computed from the
fine-tuned model, \behavior{refusal} edges above the EM control. Anchor-post
figure values appear nowhere in this computation.
